\documentclass{article}
\usepackage{fontspec}
\usepackage[a4paper,margin=1.8cm]{geometry}
\usepackage{gregoriotex}
\makeatletter
\def\gre@gregorioscore#1{%
    \begingroup%
    \let\input@path\gre@input@path%
    \input{#1.gtex}%
    \relax%
    \endgroup%
}%
\makeatother
\usepackage{pgffor}
\usepackage{needspace}
\usepackage{xcolor}

\newcommand{\scoretitle}[1]{
    \begin{center}
        \LARGE \textbf{#1}
    \end{center}
    \nopagebreak
}

\newcommand{\gregorianscore}[1]{
    \needspace{5\baselineskip}%
    \scoretitle{#1}%
    \nopagebreak
    \gregorioscore{#1}%
}

\newcommand{\gregorianpart}[2]{%
    \section*{#1}%
    \gresetgregpath{{../../../Grégorien/#1/}}%
    \foreach \gregorien in {#2} {%
        \gregorianscore{\gregorien}%
    }%
}

\begin{document}
\title{Salut du Très Saint-Sacrement}
\author{}
\date{}
\maketitle


\gregorianpart{In exhibitione}{{Adoro te},{Anima Christi},{Ave verum},{O salutaris Hostia II},{Ubi caritas}}

\gregorianpart{Ad Beatam Mariam Virginem}{{Ave Maria I},{Ave maris stella},{Inviolata}}
\subsection*{a}
\textcolor{red}{\textbf{\Vbar}}~Texte du verset \\
\textbf{\Rbar}~Texte du répons

\subsection*{d}
\textcolor{red}{\textbf{\Vbar}}~Texte du verset \\
\textbf{\Rbar}~Texte du répons

\subsection*{fe}
\textcolor{red}{\textbf{\Vbar}}~Texte du verset \\
\textbf{\Rbar}~Texte du répons

\subsection*{gr}
\textcolor{red}{\textbf{\Vbar}}~Texte du verset \\
\textbf{\Rbar}~Texte du répons


\gregorianpart{Ad Summum Pontificem}{{Tu es pastor ovium},{Tu es Petrus I},{Tu es Petrus II}}


\end{document}